% ════════════════════════════════════════════════════════════════════════════════
\section{Authentication Protocol}
\label{sec:protocol}

This section describes the protocol end-to-end, then shows how each
WireGuard parameter is produced and by which component.

\subsection{The Auth Tunnel (Noise-IK)}
\label{sec:noise}
All HTTP traffic between \texttt{UserAuthn} and the IdP passes through a
hand-built secure tunnel so the client does not need to trust a separate
TLS certificate.
We chose Noise-IK over TLS specifically because Noise-IK lets us reuse the
gateway's WireGuard public key as the responder identity, collapsing the
auth tunnel's trust anchor and the VPN data-plane trust anchor into one.

\texttt{common/noise.py} is a from-scratch implementation of the Noise-IK
handshake (Noise Protocol Framework~\cite{perrin2018noise}, revision~34) over
\texttt{cryptography.hazmat} primitives only: X25519~\cite{x25519rfc},
ChaCha20-Poly1305~\cite{rfc8439}, BLAKE2s~\cite{blake2}.
No third-party Noise library was used.

The IK pattern is a two-message handshake: in message~1 the initiator sends
its ephemeral and (encrypted) static public keys plus an encrypted payload;
in message~2 the responder replies with its ephemeral and an encrypted
payload, completing the four DH operations $es$, $ss$, $ee$, $se$.
After the handshake, both sides hold a pair of ChaCha20-Poly1305
\texttt{CipherState} objects (one per direction), and frames are exchanged
with a 4-byte big-endian length prefix (\texttt{common/framing.py}).

\subsection{End-to-End Flow}
The complete protocol proceeds in nine steps
(Fig.~\ref{fig:sequence} in Appendix~\ref{sec:appendix-sequence}):

\begin{enumerate}
  \item \texttt{UserAuthn} generates a fresh 32-byte PSK via
    \texttt{wg genpsk} on the client side.
  \item \texttt{UserAuthn} opens a TCP connection to \texttt{SrvAuthn} on
    port 9999 and runs the Noise-IK handshake. The responder static key is
    the server's WireGuard public key from the client's
    \texttt{base.conf}.
  \item Over the encrypted tunnel, \texttt{UserAuthn} sends a JSON hello
    containing service tag (\texttt{"auth"}), username, PSK, and target
    WireGuard interface name. \texttt{SrvAuthn} validates the PSK shape
    (32-byte base64) before any further processing.
  \item \texttt{SrvAuthn} calls \texttt{POST /oauth/device/code} at the
    IdP, attaching the username as a login hint and the client's source IP
    for the RBA engine.
  \item \texttt{SrvAuthn} returns \texttt{\{user\_code, verification\_uri\}}
    to the client over the tunnel; \texttt{UserAuthn} opens the URI in the
    default browser.
  \item The browser reaches the IdP through the local HTTP proxy
    (Section~\ref{sec:proxy}). The user submits the device code; the IdP
    starts the PAM-orchestrated MFA sequence (Section~\ref{sec:pam}).
  \item \texttt{SrvAuthn} polls \texttt{POST /oauth/token} every five
    seconds until the IdP returns an access token.
  \item \texttt{SrvAuthn} calls \texttt{GET /oauth/userinfo} with the
    bearer token (single-use; see section~\ref{sec:security}) and receives the
    authenticated username.
  \item \texttt{SrvAuthn} generates a fresh WG keypair
    (\texttt{wg genkey | wg pubkey}), leases the next free IP from
    \texttt{wg\_peers} inside a \texttt{BEGIN IMMEDIATE} transaction,
    registers the new peer with \texttt{wg set wg0 peer~\dots}, and sends
    the full WG config back over the tunnel.
    \texttt{UserAuthn} writes \texttt{/etc/wireguard/wg0.conf} and runs
    \texttt{wg-quick up wg0}.
\end{enumerate}


\subsection{Why OAuth Device Authorization Grant}
\texttt{UserAuthn} is a terminal application and cannot expose an HTTP
redirect URI, so a standard authorization-code flow is not an option.
The Device Authorization Grant (RFC~8628~\cite{rfc8628}) solves this by
having the terminal poll for a token while the user completes the flow in
a separate browser window, keeping the CLI and the browser completely
decoupled.

\subsection{The Local HTTP Proxy}
\label{sec:proxy}
To keep the browser flow on the same trust anchor as the auth tunnel,
\texttt{UserAuthn} listens on \texttt{127.0.0.1:18080}
(\texttt{client/proxy.py}); every browser connection opens a fresh Noise-IK
tunnel to \texttt{SrvAuthn} with the first frame
\texttt{\{"service":"proxy"\}}, and \texttt{server/proxy.py} pumps bytes
between the tunnel and the IdP at \texttt{localhost:8080}.

The IdP's \texttt{ORIGIN} is set to \texttt{http://localhost:18080}, so
every URI the IdP generates (including the OAuth
\texttt{verification\_uri\_complete}) points at the proxy.
The browser's effective trust anchor is therefore the same WireGuard
public key, just expressed through a different pipe.

WebAuthn's \texttt{RP\_ID} is \texttt{localhost}, which covers both the
proxy origin (\texttt{http://localhost:18080}) and the helpdesk's
\texttt{http://localhost:8081}, so security-key credentials registered
via the helpdesk validate at login without re-enrolment.
Using \texttt{localhost} rather than the numeric \texttt{127.0.0.1} is
required: browsers treat \texttt{localhost} as a secure context for
WebAuthn purposes, but do not consistently extend that treatment to
the IP literal.

\subsection{WireGuard Parameter Assignment}
\label{sec:wgparams}
A WireGuard peer requires six configuration values per side (client and
server).
Tables~\ref{tab:wgparams-client} and~\ref{tab:wgparams-server} show which component produces each one and how.

\begin{table}[H]
\centering
\caption{WireGuard parameters written to the client}
\label{tab:wgparams-client}
\renewcommand{\arraystretch}{1.2}
\scriptsize
\begin{tabular}{p{3.3cm}p{4.2cm}}
\toprule
\textbf{Parameter} & \textbf{How produced} \\
\midrule
1.\ Client secret key        & \texttt{wg genkey} on the server, shipped over Noise-IK \\
2.\ Client tunnel IP         & \texttt{db.lease\_next\_ip} (atomic) on the server \\
3.\ Server public key        & Pre-installed in \texttt{base.conf} \\
4.\ Server public IP         & Pre-installed in \texttt{base.conf} \\
5.\ Server UDP port          & Pre-installed in \texttt{base.conf} \\
6.\ Pre-shared key           & \texttt{wg genpsk} on the client, sent in the hello \\
\bottomrule
\end{tabular}
\end{table}

\begin{table}[H]
\centering
\caption{WireGuard parameters applied on the server}
\label{tab:wgparams-server}
\renewcommand{\arraystretch}{1.2}
\scriptsize
\begin{tabular}{p{3.3cm}p{4.2cm}}
\toprule
\textbf{Parameter} & \textbf{How produced} \\
\midrule
1.\ Server secret key        & \texttt{/etc/wireguard/privatekey-server} (generated once at setup) \\
2.\ Server tunnel IP         & \texttt{wg0.conf} \texttt{Address = 10.10.10.1/24} (written by \texttt{scripts/init-wg.sh}) \\
3.\ Server UDP port          & \texttt{srvauthn.conf} (\texttt{wg\_listen\_port}, default 51820) \\
4.\ Client public key        & \texttt{wg pubkey} on the server, then \texttt{wg set} \\
5.\ Client allowed IP        & Same lease as client param~2 \\
6.\ Pre-shared key           & Echo of client-supplied PSK \\
\bottomrule
\end{tabular}
\end{table}

The PSK (client parameter~6) is generated locally by \texttt{UserAuthn}
before the tunnel opens; the keypair and leased IP (client parameters~1
and~2) are created by \texttt{SrvAuthn} after authentication; the server's
public key, public IP, and listen port (client parameters~3--5) are
pre-distributed in \texttt{base.conf} and serve as the client's only
required prior knowledge.

One deliberate choice concerns the server tunnel IP (server parameter~2):
rather than an explicit field in \texttt{srvauthn.conf}, the server's
address (\texttt{10.10.10.1/24}) is written directly into
\texttt{wg0.conf} by \texttt{scripts/init-wg.sh} (the one-time key/config
bootstrap invoked by \texttt{scripts/setup.py}), following standard
\texttt{wg-quick} convention.
\texttt{srvauthn.conf} carries \texttt{internal\_subnet = 10.10.10.0/24},
which \texttt{SrvAuthn} uses only for client IP lease allocation (starting
from \texttt{.2}); the server permanently holds \texttt{.1}.

\subsection{Lease Lifecycle}
Every row in \texttt{wg\_peers} carries an \texttt{expires\_at}
(default 24~h, configurable in \texttt{srvauthn.conf}).
A daemon thread in \texttt{server/leases.py} polls every 60 seconds,
finds active rows with \texttt{expires\_at < now}, runs
\texttt{wg set wg0 peer <pub> remove}, and marks the row inactive.
Provisioning is the only insert path; the sweeper is the only delete path.
